\documentclass[12pt]{article}
\usepackage[T2A]{fontenc}
\usepackage[utf8]{inputenc}
\usepackage[russian]{babel}
\usepackage{amsmath,amssymb}
\usepackage{geometry}
\geometry{margin=2.5cm}
\usepackage{indentfirst}

\title{Гравитационно-волновые послесвечения от слияний\\ первичных чёрных дыр в кластеризованных гало тёмной материи}
\author{Т.~А.~Хайруллин, Ли~Чэнь}
\date{2024}

\begin{document}
\maketitle

\section*{Аннотация}
Исследуется вклад кластеризации первичных чёрных дыр (ПЧД) в темп слияний, регистрируемых LIGO/Virgo/KAGRA. Построена модель гало ПЧД с учётом динамического трения и трёхчастичных взаимодействий, предсказывающая усиление темпа слияний в маломассивном диапазоне.

\vspace{0.5cm}
\noindent\textbf{Ключевые слова:} первичные чёрные дыры; гравитационные волны; тёмная материя; темп слияний; детекторы третьего поколения.

\vspace{0.3cm}
\noindent\textbf{Раздел журнала:} Астрономия. \quad \textbf{Файл:} 1 из 3.

\section{Введение}
Первичные чёрные дыры остаются одним из кандидатов на роль тёмной материи. Кластеризация ПЧД существенно влияет на наблюдаемый темп слияний.

\section{Модель гало}
Профиль плотности гало ПЧД описывается как
\begin{equation}
\rho(r) = \frac{\rho_0}{\left(1 + r/r_s\right)^{3}},
\end{equation}
где $r_s$ — характерный радиус, $\rho_0$ — центральная плотность.

\section{Заключение}
Показано, что гравитационно-волновое послесвечение несёт информацию о профиле плотности гало.

\end{document}